% from ega2-1-fr.tex
\oldpage[II]{49}

\begin{proof}
Comme $X'$ est isomorphe à $\operatorname{Proj}(S/\mathfrak{J})$, la
condition est suffisante en vertu de (\hyperref[II.2.4.4-fr]{2.4.4}). Pour
voir qu'elle est nécessaire, considérons la suite exacte
$0\to\mathcal{J}\to\mathcal{O}_X\to\mathcal{O}_X/\mathcal{J}$, qui donne une
suite exacte
\[
  0\longrightarrow\Gamma_\bullet(\mathcal{J})
  \longrightarrow\Gamma_\bullet(\mathcal{O}_X)
  \longrightarrow\Gamma_\bullet(\mathcal{O}_X/\mathcal{J}).
\]
Il suffit de prouver que si $f\in S_m$, $g\in S_n$ sont tels que l'image dans
$\Gamma_\bullet(\mathcal{O}_X/\mathcal{J})$ de $\alpha_{n+m}(fg)$ est nulle,
alors l'une des images de $\alpha_m(f)$, $\alpha_n(g)$ est nulle. Or, par
définition, ces images sont des sections des
$(\mathcal{O}_X/\mathcal{J})$-Modules inversibles
$\mathcal{L}=(\mathcal{O}_X/\mathcal{J})(m)$ et
$\mathcal{L}'=(\mathcal{O}_X/\mathcal{J})(n)$ au-dessus du schéma intègre
$X'$~; l'hypothèse entraîne que le produit de ces deux sections est nul dans
$\mathcal{L}\otimes\mathcal{L}'$
((\hyperref[II.2.9.3-fr]{2.9.3}) et
(\hyperref[II.2.5.14.1-fr]{2.5.14.1})), donc l'une d'elles est nulle en vertu
de (\hyperref[I.7.4.4-fr]{I, 7.4.4}).
\end{proof}

\begin{corollary}[2.9.5]
\label{II.2.9.5-fr}
Soient $A$ un anneau, $M$ un $A$-module, $S$ une $A$-algèbre graduée
engendrée par l'ensemble $S_1$ des éléments homogènes de degré~$1$,
$u:M\to S_1$ un homomorphisme surjectif de $A$-modules,
$\overline{u}:\mathbf{S}(M)\to S$ l'homomorphisme (de $A$-algèbres) de
l'algèbre symétrique $\mathbf{S}(M)$ de $M$ dans $S$ qui prolonge $u$. Le
morphisme correspondant à $\overline{u}$ est alors une immersion fermée de
$\operatorname{Proj}(S)$ dans $\operatorname{Proj}(\mathbf{S}(M))$.
\end{corollary}

\begin{proof}
En effet, $\overline{u}$ est surjectif par hypothèse et il suffit d'appliquer
(\hyperref[II.2.9.2-fr]{2.9.2}).
\end{proof}

\section{Spectre homogène d'un faisceau d'algèbres graduées}
\label{section:II.3-fr}

\subsection{Spectre homogène d'une $\mathcal{O}_Y$-Algèbre graduée
quasi-cohérente.}
\label{subsection:II.3.1-fr}
\label{II.3.1-fr}

\begin{env}[3.1.1]
\label{II.3.1.1-fr}
Soient $Y$ un préschéma, $\mathcal{S}$ une $\mathcal{O}_Y$-Algèbre graduée,
$\mathcal{M}$ un $\mathcal{S}$-Module gradué. Si $\mathcal{S}$ est
\emph{quasi-cohérente}, chacun de ses composants homogènes $\mathcal{S}_n$ est
un $\mathcal{O}_Y$-Module \emph{quasi-cohérent}, étant l'image de
$\mathcal{S}$ par un homomorphisme de $\mathcal{S}$ dans lui-même
((\hyperref[I.1.3.8-fr]{I, 1.3.8}) et
(\hyperref[I.1.3.9-fr]{I, 1.3.9}))~; de même, si $\mathcal{M}$ est
quasi-cohérent en tant que $\mathcal{O}_Y$-Module, ses composants homogènes
$\mathcal{M}_n$ le sont aussi, et réciproquement. Si $d$ est un entier $>0$,
on désignera par $\mathcal{S}^{(d)}$ la somme directe des
$\mathcal{O}_Y$-Modules $\mathcal{S}_{nd}$ ($n\in\mathbf{Z}$), qui est
quasi-cohérente si $\mathcal{S}$ l'est
(\hyperref[I.1.3.9-fr]{I, 1.3.9})~; pour tout entier $k$ tel que
$0\leq k\leq d-1$, on désignera par $\mathcal{M}^{(d,k)}$ (ou
$\mathcal{M}^{(d)}$ pour $k=0$) la somme directe des
$\mathcal{M}_{nd+k}$ ($n\in\mathbf{Z}$), qui est un
$\mathcal{S}^{(d)}$-Module gradué, quasi-cohérent lorsque $\mathcal{S}$ et
$\mathcal{M}$ sont quasi-cohérents (\hyperref[I.9.6.1-fr]{I, 9.6.1}). On
notera $\mathcal{M}(n)$ le $\mathcal{S}$-Module gradué tel que
$(\mathcal{M}(n))_k=\mathcal{M}_{n+k}$ pour tout $k\in\mathbf{Z}$~; si
$\mathcal{S}$ et $\mathcal{M}$ sont quasi-cohérents, $\mathcal{M}(n)$ est un
$\mathcal{S}$-Module gradué quasi-cohérent
(\hyperref[I.9.6.1-fr]{I, 9.6.1}).

Nous dirons que $\mathcal{M}$ est un $\mathcal{S}$-Module gradué
\emph{de type fini} (resp. qu'il admet une \emph{présentation finie}) si pour
tout $y\in Y$, il existe un voisinage ouvert $U$ de $y$ et des entiers $n_i$
(resp. des entiers $m_i$ et $n_j$) tels qu'il y ait un homomorphisme surjectif
de degré~$0$
\[
  \bigoplus_{i=1}^r(\mathcal{S}(n_i)|U)\longrightarrow\mathcal{M}|U
\]
(resp. que $\mathcal{M}|U$ soit isomorphe au conoyau d'un homomorphisme de
degré~$0$
\[
  \bigoplus_{i=1}^r(\mathcal{S}(m_i)|U)
  \longrightarrow\bigoplus_{j=1}^s(\mathcal{S}(n_j)|U)).
\]

Soient $U$ un ouvert affine de $Y$, $A=\Gamma(U,\mathcal{O}_Y)$ son anneau~;
par hypothèse, la $(\mathcal{O}_Y|U)$-Algèbre graduée $\mathcal{S}|U$ est
isomorphe à $\widetilde{S}$, où $S=\Gamma(U,\mathcal{S})$ est une
$A$-algèbre

